\section{Inferencia Lógica y Demostración Matemática}

En la práctica matemática real, no construimos tablas de verdad gigantescas para validar cada paso de un razonamiento. En su lugar, utilizamos \textbf{Reglas de Inferencia}, que son esquemas lógicos válidos (tautologías probadas) que nos permiten deducir nuevas conclusiones verdaderas a partir de premisas verdaderas.

La notación estándar para un argumento consiste en listar las premisas y separarlas de la conclusión mediante una línea horizontal o el símbolo $\therefore$ (se lee ``por lo tanto'').

\subsection{Reglas de Inferencia Fundamentales}

Las siguientes son las reglas estructurales aceptadas universalmente para construir argumentos válidos:

\begin{table}[h]
\centering
\small
\renewcommand{\arraystretch}{1.5} % Más espacio vertical
\begin{tabularx}{\textwidth}{|l|c|c|X|}
\hline
\textbf{Nombre} & \textbf{Esquema Lógico} & \textbf{Tautología Base} & \textbf{Interpretación} \\ \hline
Modus Ponens & $\begin{array}{c} p \to q \\ p \\ \hline \therefore q \end{array}$ & $((p \to q) \land p) \to q$ & \textit{Afirmando el antecedente:} Si la implicación es cierta y se cumple la causa, necesariamente ocurre el efecto. \\ \hline
Modus Tollens & $\begin{array}{c} p \to q \\ \neg q \\ \hline \therefore \neg p \end{array}$ & $((p \to q) \land \neg q) \to \neg p$ & \textit{Negando el consecuente:} Si el efecto no ocurrió, la causa no pudo haber ocurrido. \\ \hline
Silogismo Hipotético & $\begin{array}{c} p \to q \\ q \to r \\ \hline \therefore p \to r \end{array}$ & $((p \to q) \land (q \to r)) \to (p \to r)$ & \textit{Transitividad:} Si A lleva a B, y B lleva a C, entonces A lleva directamente a C. \\ \hline
Silogismo Disyuntivo & $\begin{array}{c} p \lor q \\ \neg p \\ \hline \therefore q \end{array}$ & $((p \lor q) \land \neg p) \to q$ & \textit{Descarte:} Si tengo dos opciones y descarto una, debo quedarme con la otra. \\ \hline
Simplificación & $\begin{array}{c} p \land q \\ \hline \therefore p \end{array}$ & $(p \land q) \to p$ & Si sé que dos hechos son ciertos simultáneamente, puedo afirmar cualquiera de ellos por separado. \\ \hline
\end{tabularx}
\caption{Reglas de Inferencia Básicas}
\label{tab:reglas}
\end{table}

\noindent \textbf{Ejemplo de Aplicación:}
\begin{itemize}
    \item Premisa 1: Si estudio, aprobaré el examen ($p \to q$).
    \item Premisa 2: No aprobé el examen ($\neg q$).
    \item \textbf{Conclusión:} Por la regla del \textit{Modus Tollens}, puedo deducir con certeza que no estudié ($\therefore \neg p$).
\end{itemize}

\subsection{Negación de Cuantificadores}

Para demostrar la falsedad de una proposición general, es necesario saber negarla correctamente. Las Leyes de De Morgan se extienden a los cuantificadores de la siguiente manera:

\begin{enumerate}
    \item \textbf{Negación del Universal ($\forall$):}
    Para negar que ``todos'' cumplen una propiedad, no necesito probar que ninguno la cumple; basta con encontrar \textbf{al menos uno} que falle (contraejemplo).
    \[ \neg (\forall x \in A, P(x)) \iff \exists x \in A, \neg P(x) \]

    \item \textbf{Negación del Existencial ($\exists$):}
    Para negar que ``existe alguno'' que cumpla una propiedad, debo demostrar que \textbf{todos} fallan en cumplirla.
    \[ \neg (\exists x \in A, P(x)) \iff \forall x \in A, \neg P(x) \]
\end{enumerate}

\noindent \textbf{Ejemplo:}
\begin{itemize}
    \item \textit{Proposición:} ``Todos los números primos son impares'' ($\forall n \in \mathbb{P}, \text{Impar}(n)$).
    \item \textit{Negación:} ``Existe al menos un número primo que \textbf{no} es impar'' ($\exists n \in \mathbb{P}, \neg \text{Impar}(n)$).
    \item \textit{Validación:} La negación es verdadera porque existe el número $2$, que es primo y par.
\end{itemize}

\subsection{Introducción a la Demostración Formal}

Una demostración matemática es una sucesión coherente de pasos lógicos que valida la verdad de una proposición (teorema) basándose en definiciones y axiomas previos.

Para nuestros primeros ejercicios de demostración, utilizaremos las siguientes definiciones formales en el conjunto de los Enteros ($\mathbb{Z}$):

\begin{definicion}[Número Par]
Un entero $n$ es par si y solo si existe un entero $k \in \mathbb{Z}$ tal que $n = 2k$.
\end{definicion}

\begin{definicion}[Número Impar]
Un entero $n$ es impar si y solo si existe un entero $k \in \mathbb{Z}$ tal que $n = 2k + 1$.
\end{definicion}

\begin{definicion}[Divisibilidad]
Se dice que $a$ divide a $b$ (denotado $a \mid b$) si existe un entero $k \in \mathbb{Z}$ tal que $b = a \cdot k$.
\end{definicion}

\subsubsection{Demostración Modelo: Método Directo}

Analicemos paso a paso cómo se estructura una prueba formal.

\begin{teorema}
La suma de dos números enteros pares es siempre un número par.
\end{teorema}

\begin{proof}
Nuestro objetivo es probar la implicación: $a, b \text{ son pares} \implies (a+b) \text{ es par}$.

\textbf{1. Hipótesis (Datos dados):} \\
Sean $a$ y $b$ dos números enteros pares arbitrarios.

\textbf{2. Traducción usando Definiciones:} \\
Por la Definición 1 (Número Par):
\begin{itemize}
    \item Como $a$ es par, $\exists k \in \mathbb{Z}$ tal que $a = 2k$.
    \item Como $b$ es par, $\exists m \in \mathbb{Z}$ tal que $b = 2m$.
\end{itemize}

\textbf{3. Operatoria Algebraica:} \\
Sumamos ambos números:
\[ a + b = 2k + 2m \]
Factorizamos el 2 (propiedad distributiva en $\mathbb{Z}$):
\[ a + b = 2(k + m) \]

\textbf{4. Argumento Lógico (Clausura):} \\
Sabemos que la suma de dos enteros es otro entero. Sea $z = k + m$, entonces $z \in \mathbb{Z}$.
Sustituimos $z$ en la ecuación:
\[ a + b = 2z \]

\textbf{5. Conclusión:} \\
Hemos expresado $a+b$ como $2$ multiplicado por un entero. Esto satisface exactamente la definición de número par. Por lo tanto, $a+b$ es par.
\end{proof}

\subsection{Teoremas Propuestos para Demostrar}

Utilizando el esquema anterior, demuestre los siguientes teoremas:

\begin{teorema}[Producto de Par e Impar]
El producto de un entero par por cualquier entero (sea par o impar) es siempre un número par.
\end{teorema}

\begin{teorema}[Cuadrado de un Impar]
Si $n$ es un número impar, entonces $n^2$ es también impar. \\
\textit{Sugerencia: Parta de $n=2k+1$ y desarrolle el binomio cuadrado.}
\end{teorema}

\begin{teorema}[Transitividad de la Divisibilidad]
Sean $a, b, c \in \mathbb{Z}$. Si $a \mid b$ y $b \mid c$, entonces $a \mid c$.
\end{teorema}